\usepackage{hyperref}
\usepackage{tikz}
\usepackage{pbox}
\usepackage{setspace}
\usepackage{ctable}

\renewcommand*{\DefaultOptionsofRadio}{print,radio, radiosymbol=6, width=\baselineskip, bordercolor={black}, borderwidth=0pt}
\renewcommand*{\DefaultOptionsofText}{print,bordercolor={black}, backgroundcolor=white, borderwidth=0pt}
\renewcommand*{\DefaultOptionsofCheckBox}{print, borderwidth=0pt, bordercolor={0 0 0}, backgroundcolor=white}

\renewcommand{\LayoutTextField}[2]{% label, field
	\setbox0=\hbox{#1\unskip}\ifdim\wd0=0pt
	\setbox1=\hbox{#2\unskip}\ifdim\ht1>3ex
	% Multiline
	\begin{tikzpicture}[every node/.style={inner sep=0,outer sep=0}]
		\node[anchor=west] (TextFieldNode) at (0cm,0cm) {#2};
		\draw [thick] (current bounding box.south west) rectangle (current bounding box.north east);
	\end{tikzpicture}%
	\else
	% Inline field, lowered a little bit to be better integrated into the text
	\raisebox{-3.2pt}{\begin{tikzpicture}[every node/.style={inner sep=0,outer sep=0}]
			\node[anchor=west] (TextFieldNode) at (0cm,0cm) {#2};
			\draw[thick] ([yshift=-0.3ex]TextFieldNode.south west) -- ([yshift=-0.3ex]TextFieldNode.south east);
	\end{tikzpicture}}%
	\fi
	\else
	% Field with label below it
	\begin{tikzpicture}[every node/.style={inner sep=0,outer sep=0}]
		\node[anchor=west] (TextFieldNode) at (0cm,2ex) {#2};
		\draw[thick] ([yshift=-0.3ex]TextFieldNode.south west) -- ([yshift=-0.3ex]TextFieldNode.south east);
		\node[anchor=west,font=\footnotesize] at (0cm,-0.9ex) {#1};
	\end{tikzpicture}%
	\fi
}
\newcommand{\field}[2]{\TextField[width=#2]{#1}}
\newcommand{\fieldnamed}[3]{\TextField[width=#3,name=#2]{#1}}
\newcommand{\fieldinline}[2]{\TextField[width=#2,name=#1]{}}
\newcommand{\radiosize}{0.33cm}
\newcommand{\yesnoticks}[1]{%
	\raisebox{-3.2pt}{\begin{tikzpicture}[every node/.style={inner sep=0,outer sep=0}]
			\node[anchor=west,style={inner sep=2px}] (FieldYes) at (0cm,0cm) {\ChoiceMenu[radio=true,name=#1,width=\radiosize,height=\radiosize]{}{=Yes}};
			\node[anchor=west] (LabelYes) at ([xshift=0.7ex]FieldYes.east) {Yes};
			\node[anchor=west,style={inner sep=2px}] (FieldNo) at ([xshift=1ex]LabelYes.east) {\ChoiceMenu[radio=true,name=#1,width=\radiosize,height=\radiosize]{}{=No}};
			\node[anchor=west] (LabelNo) at ([xshift=0.7ex]FieldNo.east) {No};
			\draw [thick] ([xshift=-(\radiosize+0.15cm),yshift=-(\radiosize+0.15cm)]FieldYes.north east) rectangle (FieldYes.north east);
			\draw [thick] ([xshift=-(\radiosize+0.15cm),yshift=-(\radiosize+0.15cm)]FieldNo.north east) rectangle (FieldNo.north east);
	\end{tikzpicture}}%
}
\newcommand{\yesno}[2]{\pbox{0.8\textwidth}{\setstretch{1}#1}\hfill\yesnoticks{#2}}
\newcommand{\checkboxtick}[1]{%
	\raisebox{-3.2pt}{%
		\begin{tikzpicture}[every node/.style={inner sep=0,outer sep=0}]
			\node[anchor=west,inner sep=2px] (Field) at (0,0)
			{\CheckBox[name=#1,width=\radiosize,height=\radiosize]{}};
			\draw[thick]
			([xshift=-(\radiosize+0.15cm),yshift=-(\radiosize+0.15cm)]Field.north east)
			rectangle (Field.north east);
		\end{tikzpicture}%
	}%
}
\newcommand{\checkbox}[2]{\checkboxtick{#2}\hspace{0.8ex}#1}